% Pursuing Stacks (cotes 134-2 to 134-8): archive cote 134-6, site page 80 (PDF page 81), typescript page 398.
% typescript. Draft compared: ps5.tex:5529-5604.
% AI-transcribed by Claude Opus 5.5 (claude-opus-5-5, adaptive thinking, effort high), Message Batches API, batch ps3 (msgbatch_01XdFmuVVBAcgc7EFxdtCqsj), 2026-11-12; not yet checked by eye.
\scanpage{134-6/80}
\mspage{398}

(47) \[ a * G \hand{\simeq} G(a) \quad , \quad F \unsure{\hand{*}} b^{o} \hand{\simeq} F(b) \quad . \]

\uline{Remarks}.  \unsure{As|We} we are mainly interested\inserted{here} in abelianization\xout{,} and  (commutati-%
ve) homology and cohomology,  and hence in \xout{commutative} sheaves and
cosheaves with values in additive (even abelian) categories,  we are not
going to use for the time being the relationship between $A^{\wedge}$\xout{and} and $B^{\wedge}$
just touched upon.  We could elaborate a great deal more on it,  for
instance introducing a canonical pairing\inserted[typed]{(more accurately, a bi-sheaf)} with opposite variance to (42)
(48) \[ A^{\wedge o} \times B^{\wedge o} \longrightarrow (\mathrm{Sets}) \]
(or what amounts to the same,  canonical functors adjoint to each other
\[ A^{\wedge o} \to B^{\wedge} \quad , \quad B^{\wedge o} \to A^{\wedge} \quad ) \ , \]
deduced (via the \xout{isom} equivalence dual to (44)
\[ \underline{\mathrm{Hom}}^{!!}(A^{\wedge o},B^{\wedge o};\hand{\mathcal{N}}) \xrightarrow{\approx} \underline{\mathrm{Hom}}(A^{o}\times B^{o},\hand{\mathcal{N}}) \quad ) \]
from the co-pairing
\[ A^{o} \times B^{o} = B \times A \longrightarrow (\mathrm{Sets}) \]
given by
\[ (b^{o},a) \longrightarrow \mathrm{Hom}(b,a) \qquad (\text{for } a,b \ \text{ in } A\ ) \ . \]
\margin{\unsure{other description of (48) via $\times$} \unsure{cf. 397}}
The two pairings (42) and (48) can be given a common interpretation as
Hom  sets in a suitable category
\[ \hand{\mathcal{E}} = \hand{\mathcal{E}}(A) \quad , \]
which is the union of the two full subcategories  $A^{\wedge}$ and $B^{\wedge o}$  (which may
be interpreted as deduced from A by adjoining respectively small direct
and small inverse limits to it) intersecting in the common subcategory A,
the Hom sets in the two directions between an object F of $A^{\wedge}$ and an object
$G^{o}$ of $B^{\wedge o}$ (corresponding to an object G of $B^{\wedge}$ )  being given \xout{by the} respec-%
\xout{formulas} tively by the pairings (48) and (42).  The full relationship \margin{$\times$ \unsure{cf. p. 404-405}}
between these pairings is most conveniently expressed,  it seems,  by the
composition law of \xout{\unsure{maps}} maps in $\hand{\mathcal{E}}(A)$,  and associativity for this law.
The symmetry of the situation with respect to the pair (A,B)  is expressed
by the canonical isomorphism of categories
\[ \hand{\mathcal{E}}(A)^{o} \hand{\simeq} \hand{\mathcal{E}}(B) \qquad (\ \text{where } B = A^{o}\ ) \quad . \]

\inserted{\uline{Case}} \hand{(}2\textsuperscript{o}).  Of direct relevance for the abelianization story is the particu-%
lar case of the equivalence (39),  obtained by taking
\[ \hand{\mathcal{M}} = (\mathrm{Ab}) \quad . \]
Using formula(12)(page 387) for the pair $(A^{\wedge},(\mathrm{Ab}))$  and formula (14)
for $B^{\wedge}$ in guise of $\hand{\mathcal{A}}$,  we get  the canonical equivalence

\margin{80}
